\section{Methods}
\label{sec:methods}

The study evaluates a benchmark of daily PM$_{10}$ point forecasts
and audits how eligibility changes when dynamic-fidelity requirements are
added to an error-based screen. The analysis is post-evaluation: it does
not refit models or optimise a forecasting procedure.

\subsection{Data and forecasting task}

The evaluation dataset comprises daily PM$_{10}$ series from 17 Spanish
monitoring stations identified in MITECO/EMEP metadata. The stations span
urban, traffic, industrial, suburban, and rural settings. The recovered
station package does not establish one defensible common calendar endpoint
for the complete set, so we report verified station coverage and daily
resolution rather than infer a period label. The unit of analysis is a
station--model--horizon cell; the grid contains 17 stations, five model
families, and seven horizons, giving 595 structured cells. These cells
share stations and underlying series and are not independent replicates.

The task is direct multi-step daily point forecasting at horizons
$h=1,\ldots,7$. Evaluation follows five expanding rolling-origin folds in
chronological order. At each origin, information that became available
after that origin was excluded from fitting, feature construction,
preprocessing, threshold estimation, and verification inputs. Missing
verification targets were excluded rather than fabricated. Exact per-fold
row counts and one complete train/test calendar are not recoverable from
the aggregate release and are not asserted.

\subsection{Forecasting models and persistence reference}

The benchmark contains five model families: direct Histogram Gradient
Boosting (HGB), direct Ridge regression, fixed-order SARIMA, seasonal
naive, and STL+Ridge. HGB and Ridge produce direct forecasts for each
horizon. SARIMA is treated as a recursive statistical family, seasonal
naive as a seasonal carry-forward reference, and STL+Ridge as a
decomposition/regression hybrid. Full verified model-specification details
and reconstruction limits are reported in Supplementary Table S2.

Exact HGB, Ridge, and STL settings are not recoverable from the released
provenance. Historical SARIMA records conflict; we therefore do not infer
an exact order and report only the fixed-order family description supported
by the evaluation package. This limits reconstruction detail but does not
alter the aggregate diagnostics.

Persistence supplies the operational comparator. It is evaluated on the
same matched forecast cases as each candidate, so positive skill means that
the candidate reduces RMSE relative to persistence on the relevant cell.

\subsection{Evaluation metrics}

For matched forecast--observation pairs, RMSE is the square root of mean
squared forecast error. Persistence-relative skill is

\begin{equation}
\mathrm{Skill}=1-\frac{\mathrm{RMSE}_{\mathrm{model}}}
{\mathrm{RMSE}_{\mathrm{persistence}}}.
\label{eq:skill}
\end{equation}

Positive skill indicates lower RMSE than persistence; it does not by
itself establish dynamic fidelity. The primary dynamic diagnostic is the
predicted-to-observed variance ratio,

\begin{equation}
\alpha=\frac{\operatorname{Var}(\hat y)}{\operatorname{Var}(y)},
\label{eq:alpha}
\end{equation}

computed with population variances over matched cell pairs. P75 recall is
the proportion of observed P75 exceedances correctly identified by the
forecast, using a threshold estimated from the corresponding training
window.

\subsection{Eligibility rules and statistical testing}

The analysis applies two post-evaluation filters to every
station--model--horizon cell. They quantify cell eligibility; they do not
select one final winner for each station and horizon.

\begin{equation}
\textbf{Rule A:}\qquad \mathrm{Skill}>0\quad\text{and}\quad
\mathrm{DM}_{BH}\ \text{significant}.
\label{eq:rule-a}
\end{equation}

The two-sided Diebold--Mariano comparison uses the
Harvey--Leybourne--Newbold small-sample correction and horizon-dependent
autocovariance truncation. Benjamini--Hochberg adjustment is applied to
valid model--horizon p-values within each dataset before significance is
assigned at $p<0.05$.

\begin{equation}
\textbf{Rule B:}\qquad \text{Rule A}\quad\text{and}\quad
\alpha\geq0.50\quad\text{and}\quad \mathrm{Recall}_{P75}\geq0.20.
\label{eq:rule-b}
\end{equation}

The alpha and P75-recall thresholds are operational criteria for this
audit, not estimated optima, universal standards, or regulatory
recommendations. Formal error--fidelity discordance is narrower than the
geometric region in Figure~\ref{fig:skill_alpha}: a cell must pass Rule A,
have P75 recall of at least 0.20, and have $\alpha<0.50$.

\subsection{Threshold sensitivity and reproducibility}

The threshold analysis is a deterministic transformation of the analysis
rows. It changes neither forecasts nor the primary thresholds. We evaluate
the prespecified neighbourhood $\alpha\in\{0.40,0.50,0.60\}$ and
$\mathrm{Recall}_{P75}\in\{0.10,0.20,0.30\}$ to assess whether the
eligibility consequence persists; this is not a threshold-optimisation
search.

The canonical diagnostic table is the sole numerical source for the
manuscript's results and primary figures. Deterministic generators,
provenance records, and claim-level traceability are retained in the
repository package, while detailed model specification and sensitivity
tables are provided in Supplementary Material.
